\documentclass[12pt]{article}
\usepackage[utf8]{inputenc}
\usepackage[T1]{fontenc}
\usepackage{hyperref}
\title{Tzotzil, San Andr\'es Larrainzar Language Resources for K-12 Curriculum Development}
\author{Compiled by Claude Sonnet 5 (claude-sonnet-5) for Andr\'as Kornai}
\date{}
\begin{document}
\maketitle

\section{Introduction}
This report covers Tzotzil of San Andr\'es Larrainzar (ISO 639-3 tzs, Glottocode tzot1260). I found few resources for this variety. Most catalog entries file it under the wider Tzotzil code (tzo). I found no school curriculum, primers or textbooks for it in my searches.

\section{Primary Educational Resources}
I found no primers, textbooks or school programmes specific to this variety.

\section{Linguistic and Reference Materials}
\begin{itemize}
\item \emph{Diccionario Tzotzil de San Andr\'es} (1978), compiled by Alfa Hurely Vda.\ de Delgaty and Agust\'in Ruiz S\'anchez and published by SIL in Mexico. I saw only a catalog record, not the text. \url{https://glottolog.org/resource/reference/id/69837}
\end{itemize}

\section{Religious and Community Resources}
\begin{itemize}
\item Tzotzil San Andr\'es New Testament, published in 1983 by Liga B\'iblica Internacional. An electronic edition appeared in 2010. \url{https://dbs.org/bibles/TZOSAN}
\end{itemize}

\section{Digital and Audio Resources}
\begin{itemize}
\item Global Recordings Network audio in San Andr\'es Larrainzar Tzotzil, including evangelistic programs. \url{https://www.globalrecordings.net/en/language/2443}
\end{itemize}

\section{Curriculum Development Context}
I found no modern online dictionary, no corpus and no school materials for this variety. Any curriculum would have to start from the 1978 SIL dictionary and the New Testament text, and would need new material written with the community. The resources above are catalog-level findings. Their availability and condition were not checked.

\end{document}